\section{Question, background, and method}
The studied base is \texttt{\StudyModel}, at immutable revision:
\Fingerprint{\StudyRevision}
The bounded question is whether the synthetic relation can be recalled under a
fixed generation protocol without excessive transfer to close names or loss of
baseline-passing controls. The observed baseline, rather than a claim about
absence from pretraining, defines the comparison. Each attempt began with the
untouched pinned base. The manifest binds attempts to their source commits and
data hashes.\citep{historicalmanifest,retrospective}

LoRA freezes pretrained weights and learns low-rank projection updates
\citep{hu2021lora}. The complete multimodal Qwen base and native chat formatting
were retained, but inputs were text only and vision, embeddings, and the output
head remained frozen. The chosen target projections span attention,
linear-attention, and feed-forward computation. This is an implementation
boundary, not evidence that these projections uniquely store the fact.
Training used a BF16-loaded base and BF16 computation; that description alone
does not establish serialized adapter dtype.\citep{qwen35,qwen35template,retrospective}

\subsection{Prior work and adaptation}
Gangadhar and Stratos motivate conditional-target likelihood and augmentation
with unchanged facts. Their main comparisons use LoRA, while their analysis
also includes full-parameter and layer-restricted ablations. The pinned released
single-edit optimizer instead passes all model parameters to AdamW. Its loader
consumes supplied similar-fact records rather than constructing their
retrieval. The project's paper-inspired run is therefore a Qwen LoRA
adaptation, not an implementation-equivalent reproduction of either the
published experiments or that executable. The exact upstream neighbor corpus
and retrieval construction were not identified.\citep{gangadhar2024model,gangadhar2024run,gangadhar2024data}

Counterfactually augmented data motivates changing only the feature that should
change the label. Later project contrasts changed only the entity spelling;
that design was inspired by the literature but its effectiveness was not
isolated experimentally.\citep{kaushik2020cad,retrospective}

\subsection{Data, objective, and selection}
The positive-only family supervised the full target sentence. Subsequent
families supervised the object completion, adding locality or rehearsal facts
and eventually close-name contrasts. Completion-only loss masks direct loss on
prompt tokens; gradients still depend on their contextual representations.
The recipe tables distinguish configured horizons from completed steps and
selected checkpoints. The earliest family selected supervised validation loss;
the semantic family stopped at perfect behavioral validation; minimal-pair
runs completed their declared horizons and used a bounded loss tie-break after
behavioral scoring. The paper-inspired run evaluated final weights without a
validation selector.\citep{trlsft,peftlora,retrospective}

The fixed final suite never entered training or checkpoint selection. It is a
training-disjoint regression suite, not a pristine research holdout: aggregate
earlier results informed later recipes. Recall checks the taught terms,
near-name safety checks their absence for wrong names, and controls use
accepted answer aliases. Safety here does not require that every other claim
in an answer be true. Greedy, thinking-disabled generation was held fixed.
Acceptance requires all declared gates together; the source policy, not a
post-hoc score threshold, owns the decision.\citep{retrospective,historicalmanifest}

\begin{figure}[htbp]
\centering
\begin{tikzpicture}[node distance=5mm,box/.style={draw=LinkBlue,rounded corners,
  align=center,text width=0.78\linewidth,font=\small,inner sep=5pt},
  flow/.style={-{Latex},thick}]
\node[box] (base) {Pinned untouched base + hash-checked disjoint data};
\node[box,below=of base] (baseline) {Baseline generation and fixed scoring};
\node[box,below=of baseline] (train) {Frozen base + language LoRA; completion-only training};
\node[box,below=of train] (select) {Recipe-specific validation and checkpoint selection};
\node[box,below=of select] (final) {Disjoint final generation; joint acceptance decision};
\node[box,below=of final] (evidence) {Immutable evidence; later archival publication recorded separately};
\draw[flow] (base)--(baseline); \draw[flow] (baseline)--(train);
\draw[flow] (train)--(select); \draw[flow] (select)--(final);
\draw[flow] (final)--(evidence);
\end{tikzpicture}
\caption{Reviewed phase boundaries. The diagram does not assert a causal
mechanism for the learned behavior.\citep{retrospective}}
\end{figure}
